\documentclass{zhvt-classic}

% 使用本地源流明体（GenRyuMin2 TW）按文件路径加载，不依赖系统字体名
% 版框（内框／界栏／书口侧栏／鱼尾）改用与版心印章同一印泥棕色：
%   印章图片 yshs.png 实色 = RGB(135,68,52) = HTML 874434
\usepackage{xcolor}
\usepackage{tikz}
% sealbrown 双注册：xcolor 供 TikZ/pgf，L3 供 cls 版框
\definecolor{sealbrown}{HTML}{874434}
\ExplSyntaxOn
\color_set:nnn { sealbrown } { HTML } { 874434 }
\ExplSyntaxOff

\zhvtset{
  grid color       = sealbrown,
  main font        = GenRyuMin2TW-SB,
  main font path   = C:/Users/Administrator/texmf/fonts/opentype/public/GenRyuMin2TW/,
jiazhu font      = GenRyuMin2TW-M,
  jiazhu font path = C:/Users/Administrator/texmf/fonts/opentype/public/GenRyuMin2TW/,
  % 页码框容纳 4 字：正常页「九十七」+前置空 1 格=4 行恰好；
  % 100+ 页「一百零一」共 5 行超出，末字「一/二/三」被框裁掉（2026-10-04 实测），故放宽到 5
  shukou page number chars = 5,
  % 书口书名框高从默认 6 字减到 4 字，把省出的 2 字高拨给 chapter 框，
  % 使 6 字书名+卷名放进 9 字高的 chapter 框后不再压线（2026-10-05）
  shukou title chars = 4,
}

\title{史记}[淮阴侯列传]
\author{汉·司马迁}
\maker{闲叟漫记·重製}


% ===== 索隐述赞样式（2026-11-08 四次定稿 · 愚大人亲手改定参数）=====
% ⚠ 竖排前提：自定义字体家族必须带 RawFeature={vertical:+vert}；
%   simhei.ttf 加 +vhal 会整体加载失败静默回退，只写 +vert；simfang 同理。
\newCJKfontfamily{\shuzanfont}[Path=C:/Windows/Fonts/,RawFeature={vertical:+vert}]{simhei.ttf}
\newCJKfontfamily{\shuzanbody}[Path=C:/Windows/Fonts/,RawFeature={vertical:+vert}]{simfang.ttf}
% 题头：152x30pt 圆角 6pt 棕底白字；baseline (0,2) 使整框落入界行格子不压版框线；
% 白字节点 x 不补偿、y +8.2pt（像素级墨迹实测 dx/dy≈0，L/R 留白约 20pt）。
\newcommand{\shuzanhead}[1]{%
  \begin{tikzpicture}[baseline={(0,2)}]
    \fill[sealbrown, rounded corners=6pt]
      (-76pt,-15pt) rectangle (76pt,15pt);
    \node[text=white, inner sep=0pt, anchor=center] at (0pt,8.2pt)
      {\shuzanfont\fontsize{23}{23}\selectfont #1};
  \end{tikzpicture}%
}
% ★ 修正书口卷名位置：chapter 框起始位置改回 cls 默认值，避免全局影响。
%   只在每个卷首页通过 \juan 命令里的空命令微调该页位置（安全，不会跨页污染）。
\ExplSyntaxOn
\AtBeginDocument{
  \dim_set:Nn \l__zhvt_shukou_chapter_title_start_dim
    { \l__zhvt_shukou_title_height_dim + \l__zhvt_font_size_dim * 2 }
  \dim_set:Nn \l__zhvt_shukou_page_number_start_dim
    { \l__zhvt_shukou_chapter_title_start_dim + \l__zhvt_shukou_chapter_title_height_dim }
}
\ExplSyntaxOff

\begin{document}

% ★★ 版心（书口）排布覆写：把书名放到鱼尾之下（2026-10-04）
%
%   cls 默认版心自上而下：书名（title 静态框，在鱼尾之上）→ 鱼尾 → 篇名（chapter 动态框）→ 页码
%   本模板改为：          鱼尾（位置不动）→ 书名 → 篇名 → 页码分割线 → 页码
%
%   手法：把 cls 放在 title 静态框里的书名，改放进 chapter 动态框（鱼尾之下）；
%         title 静态框留空。书名取 \l__zhvt_title_str，即 \title{...} 的第一参数，
%         所以换书时只改 \title，不必在此重复写书名。
%
%   ⚠ 三条必须遵守的写法（2026-10-04 实测）：
%     ① 汉字必须用 \hbox{...} 包住。flowfram 会把框内容装进一个宽 14.4mm（= 一个 line_sep
%        = 书口栏宽）的纵向 minipage 并按该宽度折行，而整页靠 PDF /Rotate 90 显示、
%        框内文字本身不变向 —— 不包 \hbox 时「围炉夜话」4 字会被逐字折成 4 行，
%        在版面上就成了横排 4 字并横向溢出版框。
%     ② 竖排下**一个 \hbox 就是一整列**。书名与篇名要同处一列，必须放进同一个 \hbox；
%        分成两个 \hbox 时，后一个会被当成新的一列、甩到隔壁栏去
%        （实测「序」跑到书口栏外侧 14.4mm 的正文列上）。中间的 \quad 在框内是横向空白，
%        显示出来正好是书名与篇名之间的纵向间隔。
%     ③ ★ 内容盒要写成 **与框等宽的定宽盒** `\hbox to \l__zhvt_shukou_chapter_title_height_dim{…\hfil}`。
%        实测：chapter 框是把内容盒沿纵向**居中**摆放的，内容宽度一变，书名就跟着漂：
%        自然宽度的 \hbox（书名 38.4mm）落在 x=157.2mm，加篇名后（57.6mm）落在 x=152.4mm，
%        目录页（无篇名）与正文页的书名就差半格。做成等宽盒后，无论有没有篇名，
%        内容固定从框顶起排，实测两页均落在 x=142.8mm，不再漂。
\ExplSyntaxOn
% 把书口 chapter 框（书名+篇名）起始位置下移 1 字高，避免 4 字卷名顶到鱼尾
% ⚠ cls 的 \__zhvt_calculate_dimensions: 在 \AtBeginDocument 里执行，会重算该 dim；
%   必须把覆写挂到 begin document 之后，否则会被二次覆盖（2026-10-05 实测）。
% (旧方案已移到导言区，直接修改 chapter 框起始位置)
% 卷标题命令 #1=书口篇名(鱼尾) #2=正文卷名（含卷序，卷名内空格用 ~）
%
% ★★ 目录页两项修订（2026-10-05，愚大人指出）★★
%   ① 页码起点统一对齐：
%      原状 —— 目录项写成 \par\hspace{4em}卷名\hspace{4em}页码，页码起点 = 卷名实宽 + 4em。
%      卷名实宽随字数变：「卷一~梁惠王上」6 字 +~、「卷七~离娄上」5 字 +~，
%      差 1 个字位 → 卷七~卷十 4 条的页码整体比其余 10 条高 27.21pt（一个字），看着上下不齐。
%      改法 —— 卷名补全角空格（宽 = 1 字 = 27.215pt）补足到 **7 字符位**再写进目录，
%      14 条页码起点遂全部落在同一条竖线上（起点对齐；末字随页码字数自然参差）。
%      ⚠ 只对目录里的副本 \l_tmpc_tl 补，正文卷名仍用 \l_tmpa_tl 原样 —— 正文像素不受影响。
%   ② 卷名整体上移半字位：前导 \hspace 由 4em 改 3.5em。
%      实测原状目录块在版框内偏下（天头留白 139.6pt / 地脚留白 110.7pt），
%      上移 0.5em = 13.6pt 后上下基本均分，视觉居中。
\NewDocumentCommand { \juan } { mm } {
	\clearpage
	\tl_set:Nn \l__zhvt_chapter_title_tl { #1 }
	\tl_set:Nn \l_tmpa_tl { #2 }
	\tl_set_eq:NN \l_tmpc_tl \l_tmpa_tl
	\int_step_inline:nn { 7 - \str_count:n { #2 } }
	  { \tl_put_right:Nn \l_tmpc_tl { 　 } }
	\tl_set:Nx \l_tmpb_tl { \zhnum { page } }
	\tl_put_left:Nn \l_tmpb_tl { \hspace { 4em } }
	\tl_put_left:Nn \l_tmpb_tl { \l_tmpc_tl }
	\tl_put_left:Nn \l_tmpb_tl { \par \hspace { 3.5em } }
	\addtocontents { toc } { \l_tmpb_tl }
	\skip_horizontal:n { \l__zhvt_font_size_dim * 4 }
	\l_tmpa_tl
	\par
	\thispagestyle { empty }
}
\cs_set_protected:Npn \__zhvt_set_shukou_content: {
	% ★ 鱼尾上方留白（title 静态框）嵌入闲叟钤印：
	%   整页 /Rotate 90，故此处把图用 \rotatebox{90} 转正为竖章；
	%   尺寸卡在版心条幅内（书内约 7mm 宽 × 33.8mm 高，留余量不压鱼尾）。
	\setstaticcontents * { title } {
		\vbox_to_ht:nn { \l__zhvt_baseline_skip_dim } {
			\vskip 14.2pt
			\hbox_to_wd:nn { \l__zhvt_shukou_title_height_dim } {
				\hfil
				\rotatebox [ origin = c ] { 90 } { \includegraphics [ width = 7mm , height = 33.8mm ] { yshs.png } }
				\hfil
			}
		}
	}
% ★ 版心（书口）统一改法（2026-11-16 愚大人定）：18pt＋框内居中＋书口文字同版框色
%   ① 原 27.2pt 时 7 字长篇名（如「廉颇蔺相如列传」）末字溢出 chapter 框、压住页码；
%      改 \fontsize{18pt}{18pt}，前后 \hfil 令字串在框内居中（不随篇名字数漂移）。
%   ② \vbox{\vskip 8.2pt} 定量下移，抵消缩字号后栏宽方向 8.2pt 上偏（像素实测 y 中心 543.3 = 原版）。
%      ⚠ \raisebox / \vspace* 在 flowfram 框内无效；换字号须同步换补偿：16pt→10.1、19pt→7.4。
%   ③ chapter 与 page~number 均置 \color{sealbrown}（HTML 874434），与版框／界行同色；正文仍黑。
%   ⚠ 页码字位、钤印、鱼尾、界行、正文一律未动；五部合订本同此口径。
	\setdynamiccontents * { chapter }
	{ \color{sealbrown} \vbox { \vskip 8.2pt \hbox to \l__zhvt_shukou_chapter_title_height_dim
		{ \hfil \fontsize{18pt}{18pt}\selectfont \l__zhvt_title_str \quad \l__zhvt_chapter_title_tl \hfil } } }
	\setdynamiccontents * { page~number } { \color{sealbrown} \quad \zhnum { page } }
}
\__zhvt_set_shukou_content:
\ExplSyntaxOff

\maketitle

% 单篇，无多卷目录，不渲染目录页

\juan{淮阴侯列传}{卷九十二~淮阴侯列传第三十二}
\hspace*{2em}淮阴侯韩信者。淮阴人也。始为布衣时。贫。无行。不得推择为吏。又不能治生商贾。常从人寄食饮。人多厌之者。常数从其下乡南昌亭长寄食。数月。亭长妻患之。乃晨炊蓐食。食时信往。不为具食。信亦知其意。怒。竟绝去。\\
\hspace*{2em}信钓于城下。诸母漂。有一母见信饥。饭信。竟漂数十日。信喜。谓漂母曰。﹁吾必有以重报母。﹂母怒曰。﹁大丈夫不能自食。吾哀王孙而进食。岂望报乎。﹂\\
\hspace*{2em}淮阴屠中少年有侮信者。曰。﹁若虽长大。好带刀剑。中情怯耳。﹂众辱之曰。﹁信能死。刺我。不能死。出我袴下。﹂於是信孰视之。俛出袴下。蒲伏。一市人皆笑信。以为怯。\\
\hspace*{2em}及项梁渡淮。信杖剑从之。居麾下。无所知名。项梁败。又属项羽。羽以为郎中。数以策干项羽。羽不用。汉王之入蜀。信亡楚归汉。未得知名。为连敖。坐法当斩。其辈十三人皆已斩。次至信。信乃仰视。适见滕公。曰。﹁上不欲就天下乎。何为斩壮士。﹂滕公奇其言。壮其貌。释而不斩。与语。大说之。言於上。上拜以为治粟都尉。上未之奇也。信数与萧何语。何奇之。至南郑。诸将行道亡者数十人。信度何等已数言上。上不我用。即亡。何闻信亡。不及以闻。自追之。人有言上曰。﹁丞相何亡。﹂上大怒。如失左右手。居一二日。何来谒上。上且怒且喜。骂何曰。﹁若亡。何也。﹂何曰。﹁臣不敢亡也。臣追亡者。﹂上曰。﹁若所追者谁。﹂何曰。﹁韩信也。﹂上复骂曰。﹁诸将亡者以十数。公无所追。追信。诈也。﹂何曰。﹁诸将易得耳。至如信者。国士无双。王必欲长王汉中。无所事信。必欲争天下。非信无所与计事者。顾王策安所决耳。﹂王曰。﹁吾亦欲东耳。安能郁郁久居此乎。﹂何曰。﹁王计必欲东。能用信。信即留。不能用。信终亡耳。﹂王曰。﹁吾为公以为将。﹂何曰。﹁虽为将。信必不留。﹂王曰。﹁以为大将。﹂何曰。﹁幸甚。﹂於是王欲召信拜之。何曰。﹁王素慢。无礼。今拜大将如呼小兒耳。此乃信所以去也。王必欲拜之。择良日。斋戒。设坛场。具礼。乃可耳。﹂王许之。诸将皆喜。人人各自以为得大将。至拜大将。乃韩信也。一军皆惊。\\
\hspace*{2em}信拜礼毕。上坐。王曰。﹁丞相数言将军。将军何以教寡人计策。﹂信谢。因问王曰。﹁今东乡争权天下。岂非项王邪。﹂汉王曰。﹁然。﹂曰。﹁大王自料勇悍仁彊孰与项王。﹂汉王默然良久。曰。﹁不如也。﹂信再拜贺曰。﹁惟信亦为大王不如也。然臣尝事之。请言项王之为人也。项王喑噁叱咤。千人皆废。然不能任属贤将。此特匹夫之勇耳。项王见人。恭敬慈爱。言语呕呕。人有疾病。涕泣分食饮。至使人有功当封爵者。印刓敝。忍不能予。此所谓妇人之仁也。项王虽霸天下而臣诸侯。不居关中而都彭城。有背义帝之约而以亲爱王。诸侯不平。诸侯之见项王迁逐义帝。置江南。亦皆归逐其主而自王善地。项王所过。无不残灭者。天下多怨。百姓不亲附。特劫於威。彊耳。名虽为霸。实失天下心。故曰其彊易弱。今大王诚能反其道。任天下武勇。何所不诛。以天下城邑封功臣。何所不服。以义兵从思东归之士。何所不散。且三秦王为秦将。将秦子弟数岁矣。所杀亡不可胜计。又欺其众降诸侯。至新安。项王诈阬秦降卒二十余万。唯独邯。欣。翳得脱。秦父兄怨此三人。痛入骨髓。今楚彊以威王此三人。秦民莫爱也。大王之入武关。秋豪无所害。除秦苛法。与秦民约。法三章耳。秦民无不欲得大王王秦者。於诸侯之约。大王当王关中。关中民咸知之。大王失职入汉中。秦民无不恨者。今大王举而东。三秦可传檄而定也。﹂於是汉王大喜。自以为得信晚。遂听信计。部署诸将所击。\\
\hspace*{2em}八月。汉王举兵东出陈仓。定三秦。汉二年。出关。收魏。河南。韩。殷王皆降。合齐赵共击楚。四月。至彭城。汉兵败散而还。信复收兵与汉王会荥阳。复击破楚京。索之间。以故。楚兵卒不能西。\\
\hspace*{2em}汉之败却彭城。塞王欣。翟王翳亡汉降楚。齐。赵亦反汉与楚和。六月。魏王豹谒归视亲疾。至国。即绝河关反汉。与楚约和。汉王使郦生说豹。不下。其八月。以信为左丞相。击魏。魏王盛兵蒲坂。塞临晋。信乃益为疑兵。陈船欲度临晋。而伏兵从夏阳以木罂缻渡军。袭安邑。魏王豹惊。引兵迎信。信遂虏豹。定魏为河东郡。汉王遣张耳与信俱。引兵东。北击赵。代。後九月。破代兵。禽夏说阏与。信之下魏破代。汉辄使人收其精兵。诣荥阳以距楚。\\
\hspace*{2em}信与张耳以兵数万。欲东下井陉击赵。赵王。成安君陈馀闻汉且袭之也。聚兵井陉口。号称二十万。广武君李左车说成安君曰。﹁闻汉将韩信涉西河。虏魏王。禽夏说。新喋血阏与。今乃辅以张耳。议欲下赵。此乘胜而去国远斗。其锋不可当。臣闻﹃千里餽粮。士有饥色。樵苏後爨。师不宿饱﹄。今井陉之道。车不得方轨。骑不得成列。行数百里。其势粮食必在其後。愿足下假臣奇兵三万人。从间道绝其辎重。足下深沟高垒。坚营勿与战。彼前不得斗。退不得还。吾奇兵绝其後。使野无所掠。不至十日。而两将之头可致於戏下。愿君留意臣之计。否。必为二子所禽矣。﹂成安君。儒者也。常称﹁义兵不用诈谋奇计﹂。曰。﹁吾闻兵法﹃十则围之。倍则战﹄。今韩信兵号数万。其实不过数千。能千里而袭我。亦已罢极。今如此避而不击。後有大者。何以加之。则诸侯谓吾怯。而轻来伐我。﹂不听广武君策。广武君策不用。韩信使人间视。知其不用。还报。则大喜。乃敢引兵遂下。未至井陉口三十里。止舍。夜半传发。选轻骑二千人。人持一赤帜。从间道萆山而望赵军。诫曰。﹁赵见我走。必空壁逐我。若疾入赵壁。拔赵帜。立汉赤帜。﹂令其裨将传食。曰。﹁今日破赵会食。﹂诸将皆莫信。详应曰。﹁诺。﹂谓军吏曰。﹁赵已先据便地为壁。且彼未见吾大将旗鼓。未肯击前行。恐吾至阻险而还。﹂信乃使万人先行。出。背水陈。赵军望见而大笑。平旦。信建大将之旗鼓。鼓行出井陉口。赵开壁击之。大战良久。於是信。张耳详弃鼓旗。走水上军。水上军开入之。复疾战。赵果空壁争汉鼓旗。逐韩信。张耳。韩信。张耳已入水上军。军皆殊死战。不可败。信所出奇兵二千骑。共候赵空壁逐利。则驰入赵壁。皆拔赵旗。立汉赤帜二千。赵军已不胜。不能得信等。欲还归壁。壁皆汉赤帜。而大惊。以为汉皆已得赵王将矣。兵遂乱。遁走。赵将虽斩之。不能禁也。於是汉兵夹击。大破虏赵军。斩成安君泜水上。禽赵王歇。\\
\hspace*{2em}信乃令军中毋杀广武君。有能生得者购千金。於是有缚广武君而致戏下者。信乃解其缚。东乡坐。西乡对。师事之。\\
\hspace*{2em}诸将效首虏。休。毕贺。因问信曰。﹁兵法右倍山陵。前左水泽。今者将军令臣等反背水陈。曰破赵会食。臣等不服。然竟以胜。此何术也。﹂信曰。﹁此在兵法。顾诸君不察耳。兵法不曰﹃陷之死地而後生。置之亡地而後存﹄。且信非得素拊循士大夫也。此所谓﹃驱市人而战之﹄。其势非置之死地。使人人自为战。今予之生地。皆走。宁尚可得而用之乎。﹂诸将皆服曰。﹁善。非臣所及也。﹂\\
\hspace*{2em}於是信问广武君曰。﹁仆欲北攻燕。东伐齐。何若而有功。﹂广武君辞谢曰。﹁臣闻败军之将。不可以言勇。亡国之大夫。不可以图存。今臣败亡之虏。何足以权大事乎。﹂信曰。﹁仆闻之。百里奚居虞而虞亡。在秦而秦霸。非愚於虞而智於秦也。用与不用。听与不听也。诚令成安君听足下计。若信者亦已为禽矣。以不用足下。故信得侍耳。﹂因固问曰。﹁仆委心归计。原足下勿辞。﹂广武君曰。﹁臣闻智者千虑。必有一失。愚者千虑。必有一得。故曰﹃狂夫之言。圣人择焉﹄。顾恐臣计未必足用。原效愚忠。夫成安君有百战百胜之计。一旦而失之。军败鄗下。身死泜上。今将军涉西河。虏魏王。禽夏说阏与。一举而下井陉。不终朝破赵二十万众。诛成安君。名闻海内。威震天下。农夫莫不辍耕释耒。褕衣甘食。倾耳以待命者。若此。将军之所长也。然而众劳卒罢。其实难用。今将军欲举倦弊之兵。顿之燕坚城之下。欲战恐久力不能拔。情见势屈。旷日粮竭。而弱燕不服。齐必距境以自彊也。燕齐相持而不下。则刘项之权未有所分也。若此者。将军所短也。臣愚。窃以为亦过矣。故善用兵者不以短击长。而以长击短。﹂韩信曰。﹁然则何由。﹂广武君对曰。﹁方今为将军计。莫如案甲休兵。镇赵抚其孤。百里之内。牛酒日至。以飨士大夫醳兵。北首燕路。而後遣辩士奉咫尺之书。暴其所长於燕。燕必不敢不听从。燕已从。使諠言者东告齐。齐必从风而服。虽有智者。亦不知为齐计矣。如是。则天下事皆可图也。兵固有先声而後实者。此之谓也。﹂韩信曰。﹁善。﹂从其策。发使使燕。燕从风而靡。乃遣使报汉。因请立张耳为赵王。以镇抚其国。汉王许之。乃立张耳为赵王。\\
\hspace*{2em}楚数使奇兵渡河击赵。赵王耳。韩信往来救赵。因行定赵城邑。发兵诣汉。楚方急围汉王於荥阳。汉王南出。之宛。叶间。得黥布。走入成皋。楚又复急围之。六月。汉王出成皋。东渡河。独与滕公俱。从张耳军脩武。至。宿传舍。晨自称汉使。驰入赵壁。张耳。韩信未起。即其卧内上夺其印符。以麾召诸将。易置之。信。耳起。乃知汉王来。大惊。汉王夺两人军。即令张耳备守赵地。拜韩信为相国。收赵兵未发者击齐。\\
\hspace*{2em}信引兵东。未渡平原。闻汉王使郦食其已说下齐。韩信欲止。范阳辩士蒯通说信曰。﹁将军受诏击齐。而汉独发间使下齐。宁有诏止将军乎。何以得毋行也。且郦生一士。伏轼掉三寸之舌。下齐七十馀城。将军将数万众。岁馀乃下赵五十馀。为将数岁。反不如一竖儒之功乎。﹂於是信然之。从其计。遂渡河。齐已听郦生。即留纵酒。罢备汉守御。信因袭齐历下军。遂至临菑。齐王田广以郦生卖己。乃烹之。而走高密。使使之楚请救。韩信已定临菑。遂东追广至高密西。楚亦使龙且将。号称二十万。救齐。\\
\hspace*{2em}齐王广。龙且并军与信战。未合。人或说龙且曰。﹁汉兵远斗穷战。其锋不可当。齐。楚自居其地战。兵易败散。不如深壁。令齐王使其信臣招所亡城。亡城闻其王在。楚来救。必反汉。汉兵二千里客居。齐城皆反之。其势无所得食。可无战而降也。﹂龙且曰。﹁吾平生知韩信为人。易与耳。且夫救齐不战而降之。吾何功。今战而胜之。齐之半可得。何为止。﹂遂战。与信夹濰水陈。韩信乃夜令人为万馀囊。满盛沙。壅水上流。引军半渡。击龙且。详不胜。还走。龙且果喜曰。﹁固知信怯也。﹂遂追信渡水。信使人决壅囊。水大至。龙且军大半不得渡。即急击。杀龙且。龙且水东军散走。齐王广亡去。信遂追北至城阳。皆虏楚卒。\\
\hspace*{2em}汉四年。遂皆降。平齐。使人言汉王曰。﹁齐伪诈多变。反覆之国也。南边楚。不为假王以镇之。其势不定。原为假王便。﹂当是时。楚方急围汉王於荥阳。韩信使者至。发书。汉王大怒。骂曰。﹁吾困於此。旦暮望若来佐我。乃欲自立为王。﹂张良。陈平蹑汉王足。因附耳语曰。﹁汉方不利。宁能禁信之王乎。不如因而立。善遇之。使自为守。不然。变生。﹂汉王亦悟。因复骂曰。﹁大丈夫定诸侯。即为真王耳。何以假为。﹂乃遣张良往立信为齐王。徵其兵击楚。\\
\hspace*{2em}楚已亡龙且。项王恐。使盱眙人武涉往说齐王信曰。﹁天下共苦秦久矣。相与戮力击秦。秦已破。计功割地。分土而王之。以休士卒。今汉王复兴兵而东。侵人之分。夺人之地。已破三秦。引兵出关。收诸侯之兵以东击楚。其意非尽吞天下者不休。其不知厌足如是甚也。且汉王不可必。身居项王掌握中数矣。项王怜而活之。然得脱。辄倍约。复击项王。其不可亲信如此。今足下虽自以与汉王为厚交。为之尽力用兵。终为之所禽矣。足下所以得须臾至今者。以项王尚存也。当今二王之事。权在足下。足下右投则汉王胜。左投则项王胜。项王今日亡。则次取足下。足下与项王有故。何不反汉与楚连和。参分天下王之。今释此时。而自必於汉以击楚。且为智者固若此乎。﹂韩信谢曰。﹁臣事项王。官不过郎中。位不过执戟。言不听。画不用。故倍楚而归汉。汉王授我上将军印。予我数万众。解衣衣我。推食食我。言听计用。故吾得以至於此。夫人深亲信我。我倍之不祥。虽死不易。幸为信谢项王。﹂\\
\hspace*{2em}武涉已去。齐人蒯通知天下权在韩信。欲为奇策而感动之。以相人说韩信曰。﹁仆尝受相人之术。﹂韩信曰。﹁先生相人何如。﹂对曰。﹁贵贱在於骨法。忧喜在於容色。成败在於决断。以此参之。万不失一。﹂韩信曰。﹁善。先生相寡人何如。﹂对曰。﹁原少间。﹂信曰。﹁左右去矣。﹂通曰。﹁相君之面。不过封侯。又危不安。相君之背。贵乃不可言。﹂韩信曰。﹁何谓也。﹂蒯通曰。﹁天下初发难也。俊雄豪桀建号壹呼。天下之士云合雾集。鱼鳞櫜鹓。熛至风起。当此之时。忧在亡秦而已。今楚汉分争。使天下无罪之人肝胆涂地。父子暴骸骨於中野。不可胜数。楚人起彭城。转斗逐北。至於荥阳。乘利席卷。威震天下。然兵困於京。索之间。迫西山而不能进者。三年於此矣。汉王将数十万之众。距巩。雒。阻山河之险。一日数战。无尺寸之功。折北不救。败荥阳。伤成皋。遂走宛。叶之间。此所谓智勇俱困者也。夫锐气挫於险塞。而粮食竭於内府。百姓罢极怨望。容容无所倚。以臣料之。其势非天下之贤圣固不能息天下之祸。当今两主之命县於足下。足下为汉则汉胜。与楚则楚胜。臣原披腹心。输肝胆。效愚计。恐足下不能用也。诚能听臣之计。莫若两利而俱存之。参分天下。鼎足而居。其势莫敢先动。夫以足下之贤圣。有甲兵之众。据彊齐。从燕。赵。出空虚之地而制其後。因民之欲。西乡为百姓请命。则天下风走而响应矣。孰敢不听。割大弱彊。以立诸侯。诸侯已立。天下服听而归德於齐。案齐之故。有胶。泗之地。怀诸侯以德。深拱揖让。则天下之君王相率而朝於齐矣。盖闻天与弗取。反受其咎。时至不行。反受其殃。原足下孰虑之。﹂\\
\hspace*{2em}韩信曰。﹁汉王遇我甚厚。载我以其车。衣我以其衣。食我以其食。吾闻之。乘人之车者载人之患。衣人之衣者怀人之忧。食人之食者死人之事。吾岂可以乡利倍义乎。﹂蒯生曰。﹁足下自以为善汉王。欲建万世之业。臣窃以为误矣。始常山王。成安君为布衣时。相与为刎颈之交。後争张黡。陈泽之事。二人相怨。常山王背项王。奉项婴头而窜。逃归於汉王。汉王借兵而东下。杀成安君泜水之南。头足异处。卒为天下笑。此二人相与。天下至驩也。然而卒相禽者。何也。患生於多欲而人心难测也。今足下欲行忠信以交於汉王。必不能固於二君之相与也。而事多大於张黡。陈泽。故臣以为足下必汉王之不危己。亦误矣。大夫种。范蠡存亡越。霸勾践。立功成名而身死亡。野兽已尽而猎狗烹。夫以交友言之。则不如张耳之与成安君者也。以忠信言之。则不过大夫种。范蠡之於勾践也。此二人者。足以观矣。原足下深虑之。且臣闻勇略震主者身危。而功盖天下者不赏。臣请言大王功略。足下涉西河。虏魏王。禽夏说。引兵下井陉。诛成安君。徇赵。胁燕。定齐。南摧楚人之兵二十万。东杀龙且。西乡以报。此所谓功无二於天下。而略不世出者也。今足下戴震主之威。挟不赏之功。归楚。楚人不信。归汉。汉人震恐。足下欲持是安归乎。夫势在人臣之位而有震主之威。名高天下。窃为足下危之。﹂韩信谢曰。﹁先生且休矣。吾将念之。﹂\\
\hspace*{2em}後数日。蒯通复说曰。﹁夫听者事之候也。计者事之机也。听过计失而能久安者。鲜矣。听不失一二者。不可乱以言。计不失本末者。不可纷以辞。夫随厮养之役者。失万乘之权。守儋石之禄者。阙卿相之位。故知者决之断也。疑者事之害也。审豪氂之小计。遗天下之大数。智诚知之。决弗敢行者。百事之祸也。故曰﹃猛虎之犹豫。不若蜂虿之致螫。骐骥之跼躅。不如驽马之安步。孟贲之狐疑。不如庸夫之必至也。虽有舜禹之智。吟而不言。不如瘖聋之指麾也﹄。此言贵能行之。夫功者难成而易败。时者难得而易失也。时乎时。不再来。原足下详察之。﹂韩信犹豫不忍倍汉。又自以为功多。汉终不夺我齐。遂谢蒯通。蒯通说不听。已详狂为巫。\\
\hspace*{2em}汉王之困固陵。用张良计。召齐王信。遂将兵会垓下。项羽已破。高祖袭夺齐王军。汉五年正月。徙齐王信为楚王。都下邳。\\
\hspace*{2em}信至国。召所从食漂母。赐千金。及下乡南昌亭长。赐百钱。曰。﹁公。小人也。为德不卒。﹂召辱己之少年令出胯下者以为楚中尉。告诸将相曰。﹁此壮士也。方辱我时。我宁不能杀之邪。杀之无名。故忍而就於此。﹂\\
\hspace*{2em}项王亡将锺离眜家在伊庐。素与信善。项王死后。亡归信。汉王怨眜。闻其在楚。诏楚捕眜。信初之国。行县邑。陈兵出入。汉六年。人有上书告楚王信反。高帝以陈平计。天子巡狩会诸侯。南方有云梦。发使告诸侯会陈。﹁吾将游云梦。﹂实欲袭信。信弗知。高祖且至楚。信欲发兵反。自度无罪。欲谒上。恐见禽。人或说信曰。﹁斩眜谒上。上必喜。无患。﹂信见眜计事。眜曰。﹁汉所以不击取楚。以眜在公所。若欲捕我以自媚於汉。吾今日死。公亦随手亡矣。﹂乃骂信曰。﹁公非长者。﹂卒自刭。信持其首。谒高祖於陈。上令武士缚信。载后车。信曰。﹁果若人言。﹃狡兔死。良狗烹。高鸟尽。良弓藏。敌国破。谋臣亡。﹄天下已定。我固当烹。﹂上曰。﹁人告公反。﹂遂械系信。至雒阳。赦信罪。以为淮阴侯。\\
\hspace*{2em}信知汉王畏恶其能。常称病不朝从。信由此日夜怨望。居常鞅鞅。羞与绛。灌等列。信尝过樊将军哙。哙跪拜送迎。言称臣。曰。﹁大王乃肯临臣。﹂信出门。笑曰。﹁生乃与哙等为伍。﹂上常从容与信言诸将能不。各有差。上问曰。﹁如我能将几何。﹂信曰。﹁陛下不过能将十万。﹂上曰。﹁於君何如。﹂曰。﹁臣多多而益善耳。﹂上笑曰。﹁多多益善。何为为我禽。﹂信曰。﹁陛下不能将兵。而善将将。此乃信之所以为陛下禽也。且陛下所谓天授。非人力也。﹂\\
\hspace*{2em}陈豨拜为钜鹿守。辞於淮阴侯。淮阴侯挈其手。辟左右与之步於庭。仰天叹曰。﹁子可与言乎。欲与子有言也。﹂豨曰。﹁唯将军令之。﹂淮阴侯曰。﹁公之所居。天下精兵处也。而公。陛下之信幸臣也。人言公之畔。陛下必不信。再至。陛下乃疑矣。三至。必怒而自将。吾为公从中起。天下可图也。﹂陈豨素知其能也。信之。曰。﹁谨奉教。﹂汉十年。陈豨果反。上自将而往。信病不从。阴使人至豨所。曰。﹁弟举兵。吾从此助公。﹂信乃谋与家臣夜诈诏赦诸官徒奴。欲发以袭吕后。太子。部署已定。待豨报。其舍人得罪於信。信囚。欲杀之。舍人弟上变。告信欲反状於吕后。吕后欲召。恐其党不就。乃与萧相国谋。诈令人从上所来。言豨已得死。列侯群臣皆贺。相国绐信曰。﹁虽疾。彊入贺。﹂信入。吕后使武士缚信。斩之长乐锺室。信方斩。曰。﹁吾悔不用蒯通之计。乃为兒女子所诈。岂非天哉。﹂遂夷信三族。\\
\hspace*{2em}高祖已从豨军来。至。见信死。且喜且怜之。问。﹁信死亦何言。﹂吕后曰。﹁信言恨不用蒯通计。﹂高祖曰。﹁是齐辩士也。﹂乃诏齐捕蒯通。蒯通至。上曰。﹁若教淮阴侯反乎。﹂对曰。﹁然。臣固教之。竖子不用臣之策。故令自夷於此。如彼竖子用臣之计。陛下安得而夷之乎。﹂上怒曰。﹁烹之。﹂通曰。﹁嗟乎。冤哉烹也。﹂上曰。﹁若教韩信反。何冤。﹂对曰。﹁秦之纲绝而维弛。山东大扰。异姓并起。英俊乌集。秦失其鹿。天下共逐之。於是高材疾足者先得焉。蹠之狗吠尧。尧非不仁。狗因吠非其主。当是时。臣唯独知韩信。非知陛下也。且天下锐精持锋欲为陛下所为者甚众。顾力不能耳。又可尽烹之邪。﹂高帝曰。﹁置之。﹂乃释通之罪。\\
\hspace*{2em}太史公曰。吾入淮阴。淮阴人为余言。韩信虽为布衣时。其志与众异。其母死。贫无以葬。然乃行营高敞地。令其旁可置万家。余视其母冢。良然。假令韩信学道谦让。不伐己功。不矜其能。则庶几哉。於汉家勋可以比周。召。太公之徒。後世血食矣。不务出此。而天下已集。乃谋畔逆。夷灭宗族。不亦宜乎。\\

% ===== 索隐述赞（司马贞《史记索隐》）· 统一样式 =====
% 题头「索隐述赞曰」五字并入反白框；正文每四句一列，定宽盒 \hfil 居列中；
% \fontsize{24}{41} 列距对齐全书界行（121+41k）；\clearpage 使述赞独占一页。
\clearpage
\shuzanhead{索隐述赞曰}\linebreak
{\shuzanbody\fontsize{24}{41}\selectfont
\hbox to \hsize{\hfil 君臣一体。自古所难。相国深荐。策拜登坛。\hfil}\linebreak
\hbox to \hsize{\hfil 沈沙决水。拔帜传餐。与汉汉重。归楚楚安。\hfil}\linebreak
\hbox to \hsize{\hfil 三分不议。伪游可叹。\hfil}%
\par}
\end{document}
